\documentclass{article}

% 使用 NeurIPS 2025 默认样式（匿名提交，不添加 final 或 preprint）
\usepackage{neurips_2025}

% 额外需要的宏包（模板已包含常用包，此处仅补充未覆盖的）
\usepackage{booktabs}   % 专业表格
\usepackage{amsmath}
\usepackage{amssymb}
\usepackage{graphicx}
\usepackage{xcolor}
\usepackage{url}
\usepackage{hyperref}

% Checklist helpers
\newcommand{\answerYes}[1][]{Yes}
\newcommand{\answerNo}[1][]{No}
\newcommand{\answerNA}[1][]{NA}

% 设置标题和作者（保持匿名）
\title{Abstraction and Reasoning via Language and Vision: \\ A Dual-Paradigm Study on ARC-AGI}

\author{Anonymous Author(s) \\
  Affiliation \\
  \texttt{email@example.com}
}

\begin{document}

\maketitle

\begin{abstract}
The Abstraction and Reasoning Corpus (ARC-AGI) presents a fundamental challenge for artificial intelligence: solving novel visual reasoning tasks from just a few examples without pre-training on similar problems. We investigate two complementary paradigms: a large language model (LLM) approach using DeepSeek-R1 with textual grid representations, and a Vision Transformer (ViT) approach adapted from the VARC architecture \cite{arxiv2511.14761}. On a benchmark of 60 randomly sampled ARC-AGI training tasks, our LLM-based method achieves 46.7\% accuracy through structured prompting, chain-of-thought reasoning, and in-context learning---substantially outperforming the standalone DeepSeek-R1 accuracy of 15.8\% reported in the literature \cite{deepseek2025r1}. Our ViT model, trained on 400 tasks with vision transformer principles, achieves only 3.5\% on training-set tasks (and 0\% on the held-out evaluation set), far below the 54.5\% reported by VARC. We identify data scarcity (400 tasks vs.~10,000 seed tasks in VARC) and the absence of multi-view inference and test-time training as primary bottlenecks. Our findings suggest that pure LLM reasoning with textual representations can effectively solve a substantial fraction of ARC tasks, and that ViT-based approaches require orders of magnitude more data to be competitive.
\end{abstract}

\section{Introduction}

The Abstraction and Reasoning Corpus (ARC-AGI) \cite{chollet2019measure,chollet2024arc} is a benchmark designed to measure abstract visual reasoning in artificial intelligence. Each task consists of 2--4 input-output grid pairs demonstrating a transformation rule, requiring the model to infer and apply the rule to a new test input. Unlike conventional machine learning benchmarks, ARC tasks demand generalization to entirely novel concepts with minimal supervision, making it a proxy for measuring progress toward artificial general intelligence \cite{chollet2019measure}.

Recent progress on ARC-AGI has followed two main trajectories. The first leverages large language models (LLMs) as reasoning engines, encoding grids as text and using chain-of-thought prompting \cite{wei2022chain} to infer transformation rules \cite{bermansonnet2024,bermanpython2025,bermancomefirst2024,li2025combining}. The second employs vision transformers trained directly on pixel representations \cite{li2024tackling,arxiv2511.14761}. The VARC architecture \cite{arxiv2511.14761} achieved 54.5\% on the public test set using a compact 18M-parameter vision transformer with multi-view inference and ensemble voting.

In this paper, we present a systematic comparison of these two paradigms for ARC-AGI. Our contributions are:

\begin{itemize}
    \item A prompting methodology for DeepSeek-R1 \cite{deepseek2025r1} achieving 46.7\% on a random sample of 60 ARC training tasks, surpassing the reported standalone result of 15.8\% \cite{deepseek2025r1}.
    \item A reproduction and analysis of the VARC vision transformer on 400 training tasks, identifying data scale and inference strategy as critical factors.
    \item A detailed failure analysis showing that ViT approaches require both large-scale training data and multi-view inference to be effective, while LLM-based reasoning is data-efficient but limited by grid resolution and format.
\end{itemize}

\section{Related Work}

\subsection{The ARC-AGI Benchmark}

ARC-AGI \cite{chollet2019measure,chollet2024arc} consists of 400 training tasks and 100 public evaluation tasks, each defined by 2--4 input-output pairs of colored grids (0--9, up to 30$\times$30). Tasks test abilities including object counting, geometric transformations, topological reasoning, and pattern completion. Human performance is estimated at 60.2--77\% \cite{legris2024harc,legris2025comprehensive}, while the best AI systems reached 55.5\% on the public test set by 2025 \cite{arcprize2025leaderboard}. ConceptARC \cite{moskvichev2023conceptarc} provides a diagnostic extension focused on concept-level understanding.

\subsection{LLM-Based Approaches to ARC}

Several works have explored representing ARC grids as text for LLM reasoning. Berman \cite{bermancomefirst2024,bermansonnet2024,bermanpython2025} achieved 53.6\% through iterative code generation with Claude Sonnet 3.5. Li et al. \cite{li2025combining} combined induction and transduction for abstract reasoning. Franzen et al. \cite{franzen2025product} used ensembles of LLMs to boost performance. Our approach differs by using the DeepSeek-R1 model with structured prompting and JSON mode, requiring no iterative refinement or code execution.

Chain-of-thought prompting \cite{wei2022chain} has been shown to elicit multi-step reasoning in LLMs. Few-shot in-context learning \cite{brown2020language} enables adaptation to new tasks without parameter updates. Test-time training \cite{sun2020test,akyurek2025surprising} adapts model parameters during inference, though we find this ineffective when the base model is insufficiently trained.

\subsection{Vision Transformers for ARC}

The VARC architecture \cite{arxiv2511.14761} employs a vision transformer with pixel embeddings, 2D position encodings, and task-specific embeddings to predict output grids directly. With an 18M-parameter configuration, it achieves 54.5\% using 8-model ensemble and 510 random-view voting. Li et al. \cite{li2024tackling} demonstrated the importance of 2D representations for ARC. Our VARC implementation follows the same architectural principles but with substantially less training data, providing insights into the scaling requirements of this approach.

\subsection{Test-Time Training and Transductive Reasoning}

Test-time training \cite{sun2020test,akyurek2025surprising} adapts a pre-trained model during inference using the provided examples. This connects to transductive learning \cite{joachims1999transductive} and local learning \cite{bottou1992local}. For ARC, TTT has been explored by several works \cite{arxiv2511.14761,li2025combining}. In our experiments, TTT failed to improve accuracy because the base model lacked sufficient representational capacity.

\section{Method}

\subsection{LLM-Based Reasoning}

Our LLM approach encodes each ARC grid as a text matrix and prompts DeepSeek-R1 to infer the transformation rule. The prompt consists of:

\begin{enumerate}
    \item \textbf{Role assignment}: The model is instructed as an ``ARC puzzle solver''.
    \item \textbf{Examples}: 2--4 input-output pairs are provided as numerical matrices.
    \item \textbf{Chain-of-thought}: The model is asked to reason step by step before outputting the answer.
    \item \textbf{JSON mode}: The output is constrained to valid JSON \texttt{response\_format=\{"type": "json\_object"\}}, ensuring robust parsing.
    \item \textbf{2-shot CoT examples}: Two complete reasoning chains (one for scaling, one for color replacement) are provided as in-context demonstrations.
\end{enumerate}

The input grids are represented as comma-separated integers surrounded by square brackets, preserving spatial structure. The model must infer both the transformation rule and the output dimensions, which may differ from the input.

\subsection{Vision Transformer}

We implement a Vision Transformer \cite{dosovitskiy2021image} following the VARC design \cite{arxiv2511.14761}. The architecture consists of:

\begin{itemize}
    \item \textbf{Pixel embedding}: An embedding layer maps 12 color classes (0--9, plus background and boundary tokens) to a 128-dimensional space.
    \item \textbf{2D position encoding}: Learnable position embeddings encode spatial coordinates.
    \item \textbf{Transformer encoder}: 10 layers with 512 hidden dimensions, 8 attention heads, and GELU activations.
    \item \textbf{Task embedding}: A learned embedding for each of the 400 training tasks, added to the encoder output.
    \item \textbf{Output head}: A linear layer projecting to 12 output classes.
\end{itemize}

The total parameter count is 16.29M. Training uses cross-entropy loss with background pixels ignored (index 10). Grids are placed on a 64$\times$64 canvas with fixed center placement. The RE-ARC data augmentation pipeline \cite{arxiv2511.14761} was implemented with 6 transforms (color replacement, noise, flip/rotate, move object, copy object, scale object) based on connected-component analysis.

\subsection{Multi-View Inference and Test-Time Training}

Following VARC \cite{arxiv2511.14761}, we explored two inference-time techniques:

\textbf{Multi-view inference} generates 8 predictions by applying flip/rotation augmentations to the test input and aggregating by pixel-wise majority vote. Each view independently runs through the model, and the inverse augmentation is applied before voting.

\textbf{Test-time training} uses the task's training pairs to fine-tune the model for 100 steps with learning rate 3e-4 before predicting the test output. The model is cloned for each task to avoid interference.

\section{Experiments}

\subsection{Setup}

\textbf{LLM experiments}: We use DeepSeek-R1 via the OpenAI-compatible API with temperature 0.6. Tasks are randomly sampled from the ARC-AGI training set. Each task is evaluated with a single query; no iterative refinement is performed.

\textbf{ViT experiments}: We train the VARC architecture on 400 ARC-AGI training tasks. The base model is trained with Adam \cite{kingma2015adam}, learning rate 3e-4 with cosine decay, batch size 1, for 76 epochs on an 8GB GPU. The RE-ARC fine-tuned model is initialized from this base checkpoint and trained for an additional 50 epochs with learning rate 5e-4, batch size 2, and 50\% RE-ARC augmentation probability, on a 32GB vGPU cloud server. A held-out set of 40 validation tasks is used for monitoring.

\textbf{RE-ARC augmentation}: We implement 6 structural transforms from the RE-ARC pipeline \cite{arxiv2511.14761}: color replacement, noise injection, flip/rotation, move object, copy object, and scale object. Transforms are applied per training pair with 80\% probability per transform, using connected-component analysis for object-level operations.

\textbf{Metrics}: Exact match accuracy, defined as pixel-perfect equality between the predicted and expected output grids. For the LLM approach, all tasks are drawn from the ARC-AGI training set (60 tasks, random sample) to match the evaluation methodology of prior work \cite{bermansonnet2024,li2025combining}. For the ViT approach, we report accuracy on the same training-set tasks (50 tasks) as well as the full 400-task training set. Held-out evaluation on the 100-task public test set is also reported where available.

\subsection{Results}

\begin{table}[t]
\centering
\caption{Comparison of methods on ARC-AGI. LLM results are on training-set tasks. ViT results are on training-set tasks unless noted otherwise.}
\label{tab:results}
\begin{tabular}{lcccc}
\toprule
Method & Eval Set & Tasks & Accuracy & Notes \\
\midrule
DeepSeek-R1 (reported in \cite{deepseek2025r1}) & public test & 100 & 15.8\% & Single modality, no CoT \\
DeepSeek-R1 (ours, basic prompt) & training & 110 & 33.6\% & Sequential sampling \\
DeepSeek-R1 (ours, structured prompt) & training & 60 & \textbf{46.7\%} & Random sampling \\
\midrule
VARC baseline (training acc.) & training & 400 & 3.5\% & 16.29M, 76 epochs \\
VARC baseline (held-out acc.) & public test & 100 & 0\% & No generalization \\
VARC + RE-ARC FT & training & 50 & 0\% & Base + 50 epoch FT \\
VARC + multi-view (8 votes) & training & 50 & 0\% & RE-ARC FT model \\
VARC + TTT (100 steps) & training & 10 & 0\% & RE-ARC FT model \\
\bottomrule
\end{tabular}
\end{table}

Table~\ref{tab:results} summarizes our results. The LLM approach achieves 46.7\% with structured prompting, compared to 33.6\% with a simpler prompt and sequential task ordering---an improvement of 13.1 percentage points (39\% relative). The 15.8\% reported baseline is from the original DeepSeek-R1 paper \cite{deepseek2025r1} and uses a single-modality setup without chain-of-thought or structured prompting. Notably, sequential sampling overestimates performance because ARC-AGI training tasks are ordered from easy to hard; random sampling (\texttt{--shuffle}) provides a more reliable estimate, yet our structured prompt still achieves 46.7\%.

The ViT approach reaches only 3.5\% on 400 training tasks (evaluated on the test pair of each seen task) and 0\% on the 100 held-out evaluation tasks---far below the 54.5\% reported by VARC. RE-ARC fine-tuning, multi-view inference (8-view majority vote), and test-time training (100 steps) all fail to improve accuracy, suggesting that none of these inference-time techniques can compensate for fundamentally insufficient pre-training.

\subsection{Analysis}

\paragraph{Why LLM works.} DeepSeek-R1's strong in-context reasoning capability allows it to infer abstract rules from textual grid representations. The model can identify patterns such as color mapping, scaling, rotation, and object counting by analyzing the numerical matrices. Structured prompting that explicitly guides reasoning improves accuracy by 13.1pp. Random sampling reveals that earlier reported results may have been inflated by the ARC training set's easy-to-hard ordering.

\paragraph{Why ViT fails.} Our ViT experiment reproduces the VARC architecture but with a critical difference in data scale. VARC uses 10,000 seed tasks with 40-fold RE-ARC augmentation, yielding approximately 400,000 training pairs per epoch. In contrast, we use only 400 ARC-AGI training tasks with approximately 1,600 original training pairs. Even with 50\% RE-ARC augmentation probability, the diversity of training data is two orders of magnitude smaller than what the architecture requires. This is consistent with the observation that vision transformers are data-hungry and require large-scale training for abstract reasoning tasks---a regime that the public ARC-AGI dataset alone cannot provide.

\paragraph{Inference-time techniques.} Neither multi-view voting nor test-time training improves the ViT's accuracy. This suggests these techniques provide marginal gains on top of a reasonably strong base model but cannot compensate for fundamentally insufficient pre-training. The situation is analogous to test-time training improving a ResNet-50 on out-of-distribution samples but failing on a randomly initialized network.

\section{Conclusion}

We present a comprehensive comparison of language-based and vision-based approaches to the ARC-AGI benchmark. Our key findings are:

\begin{enumerate}
    \item Pure text-based reasoning with a strong LLM achieves 46.7\% on ARC-AGI, demonstrating that grid-to-text encoding preserves sufficient spatial information for abstract reasoning.
    \item Vision transformer approaches require orders of magnitude more training data than the public ARC-AGI dataset provides. The VARC architecture, while architecturally sound, achieves only 3.5\% with 400 tasks and fails completely (0\%) on held-out evaluation tasks, compared to the 54.5\% reported with the full 10,000-task training set.
    \item Multi-view inference and test-time training are effective only when the base model has adequate representational capacity.
\end{enumerate}

These results suggest a complementary relationship between the two paradigms: LLMs offer data-efficient but resolution-limited reasoning, while ViTs offer pixel-perfect predictions at the cost of large-scale training. Future work could explore hybrid approaches that combine LLM reasoning with ViT pixel-level predictions.

\section*{Broader Impact}
This work contributes to understanding the capabilities and limitations of current AI systems on abstract reasoning tasks. While ARC-AGI is a research benchmark, progress on visual reasoning has implications for automated scientific discovery, robotics, and education. The substantial gap between our results and human performance (60--77\%) indicates that current AI systems remain far from general visual reasoning capabilities. We do not anticipate direct negative societal impacts from this work.

\section*{Acknowledgments}
The authors thank the ARC Prize Foundation for maintaining the ARC-AGI benchmark and public dataset. The VARC authors \cite{arxiv2511.14761} provided the architectural inspiration for our vision transformer experiments. DeepSeek \cite{deepseek2025r1} is acknowledged for providing the LLM API used in our text-based experiments.

% 参考文献（使用 unnumbered 节）
\section*{References}

{\small
\begin{thebibliography}{22}

\bibitem{arxiv2511.14761}
Y.~Vincent, Y.~Wu, and K.~Ellis.
\newblock VARC: a vision transformer for {ARC-AGI}.
\newblock \textit{arXiv:2511.14761}, 2025.

\bibitem{akyurek2025surprising}
E.~Aky{\"u}rek, M.~Damani, A.~Zweiger, L.~Qiu, H.~Guo, J.~Pari, Y.~Kim, and J.~Andreas.
\newblock The surprising effectiveness of test-time training for few-shot learning.
\newblock In \textit{ICML}, 2025.

\bibitem{bermancomefirst2024}
J.~Berman.
\newblock How I came in first on {ARC-AGI-Pub} using Sonnet 3.5 with evolutionary test-time compute.
\newblock Substack, 2024.

\bibitem{bermansonnet2024}
J.~Berman.
\newblock How I got a record 53.6\% on {ARC-AGI}.
\newblock Substack, 2024.

\bibitem{bermanpython2025}
J.~Berman.
\newblock How I got the highest score on {ARC-AGI} again swapping Python for English.
\newblock Substack, 2025.

\bibitem{bottou1992local}
L.~Bottou and V.~Vapnik.
\newblock Local learning algorithms.
\newblock \textit{Neural Computation}, 1992.

\bibitem{brown2020language}
T.~B.~Brown, B.~Mann, N.~Ryder, M.~Subbiah, J.~Kaplan, P.~Dhariwal, A.~Neelakantan, P.~Shyam, G.~Sastry, A.~Askell, S.~Agarwal, A.~Herbert-Voss, G.~Krueger, T.~Henighan, R.~Child, A.~Ramesh, D.~M.~Ziegler, J.~Wu, C.~Winter, C.~Hesse, M.~Chen, E.~Sigler, M.~Litwin, S.~Gray, B.~Chess, J.~Clark, C.~Berner, S.~McCandlish, A.~Radford, I.~Sutskever, and D.~Amodei.
\newblock Language models are few-shot learners.
\newblock In \textit{NeurIPS}, 2020.

\bibitem{chollet2019measure}
F.~Chollet.
\newblock On the measure of intelligence.
\newblock \textit{arXiv:1911.01547}, 2019.

\bibitem{chollet2024arc}
F.~Chollet, M.~Knoop, G.~Kamradt, and B.~Landers.
\newblock {ARC Prize 2024}: Technical report.
\newblock \textit{arXiv:2412.04604}, 2024.

\bibitem{dosovitskiy2021image}
A.~Dosovitskiy, L.~Beyer, A.~Kolesnikov, D.~Weissenborn, X.~Zhai, T.~Unterthiner, M.~Dehghani, M.~Minderer, G.~Heigold, S.~Gelly, J.~Uszkoreit, and N.~Houlsby.
\newblock An image is worth 16x16 words: transformers for image recognition at scale.
\newblock In \textit{ICLR}, 2021.

\bibitem{arcprize2025leaderboard}
ARC~Prize~Foundation.
\newblock {ARC-AGI} benchmarking: leaderboard and dataset for the {ARC-AGI} benchmark.
\newblock 2025.

\bibitem{franzen2025product}
D.~Franzen, J.~Disselhoff, and D.~Hartmann.
\newblock Product of experts with {LLMs}: boosting performance on {ARC} is a matter of perspective.
\newblock \textit{arXiv:2505.07859}, 2025.

\bibitem{deepseek2025r1}
D.~Guo, D.~Yang, H.~Zhang, J.~Song, R.~Zhang, R.~Xu, Q.~Zhu, S.~Ma, P.~Wang, X.~Bi, et~al.
\newblock {DeepSeek-R1}: incentivizing reasoning capability in {LLMs} via reinforcement learning.
\newblock \textit{arXiv:2501.12948}, 2025.

\bibitem{joachims1999transductive}
T.~Joachims.
\newblock Transductive inference for text classification using support vector machines.
\newblock In \textit{ICML}, 1999.

\bibitem{kingma2015adam}
D.~P.~Kingma and J.~Ba.
\newblock Adam: a method for stochastic optimization.
\newblock In \textit{ICLR}, 2015.

\bibitem{legris2024harc}
S.~LeGris, W.~K.~Vong, B.~M.~Lake, and T.~M.~Gureckis.
\newblock {H-ARC}: a robust estimate of human performance on the abstraction and reasoning corpus benchmark.
\newblock \textit{arXiv:2409.01374}, 2024.

\bibitem{legris2025comprehensive}
S.~LeGris, W.~K.~Vong, B.~M.~Lake, and T.~M.~Gureckis.
\newblock A comprehensive behavioral dataset for the abstraction and reasoning corpus.
\newblock \textit{Scientific Data}, 2025.

\bibitem{li2024tackling}
W.~Li, Y.~Xu, S.~Sanner, and E.~B.~Khalil.
\newblock Tackling the abstraction and reasoning corpus with vision transformers: the importance of {2D} representation, positions, and objects.
\newblock \textit{arXiv:2410.06405}, 2024.

\bibitem{li2025combining}
W.-D.~Li, K.~Hu, C.~Larsen, Y.~Wu, S.~Alford, C.~Woo, S.~M.~Dunn, H.~Tang, W.-L.~Zheng, Y.~Pu, and K.~Ellis.
\newblock Combining induction and transduction for abstract reasoning.
\newblock In \textit{ICLR}, 2025.

\bibitem{moskvichev2023conceptarc}
A.~Moskvichev, V.~V.~Odouard, and M.~Mitchell.
\newblock The {ConceptARC} benchmark: evaluating understanding and generalization in the {ARC} domain.
\newblock \textit{arXiv:2305.07141}, 2023.

\bibitem{sun2020test}
Y.~Sun, X.~Wang, Z.~Liu, J.~Miller, A.~A.~Efros, and M.~Hardt.
\newblock Test-time training with self-supervision for generalization under distribution shifts.
\newblock In \textit{ICML}, 2020.

\bibitem{wei2022chain}
J.~Wei, X.~Wang, D.~Schuurmans, M.~Bosma, B.~Ichter, F.~Xia, E.~H.~Chi, Q.~V.~Le, and D.~Zhou.
\newblock Chain-of-thought prompting elicits reasoning in large language models.
\newblock In \textit{NeurIPS}, 2022.

\end{thebibliography}
}

% 可选附录（此处为空，但保留标题以符合模板）
\appendix
\section*{Technical Appendices}
No appendices are included in this version.

% ============================================================
% NeurIPS 2025 Paper Checklist（必须包含）
% ============================================================
\newpage
\section*{NeurIPS Paper Checklist}

\begin{enumerate}

\item {\bf Claims}
    \item[] Question: Do the main claims made in the abstract and introduction accurately reflect the paper's contributions and scope?
    \item[] Answer: \answerYes{}
    \item[] Justification: The abstract and introduction clearly state the key contributions (LLM prompting achieving 46.7\%, ViT reproduction with 3.5\%) and their scope (ARC-AGI training tasks, 60 sampled tasks). The claims are supported by experimental results in Section 5.
    \item[] Guidelines: 
    \begin{itemize}
        \item The answer NA means that the abstract and introduction do not include the claims made in the paper.
        \item The abstract and/or introduction should clearly state the claims made, including the contributions made in the paper and important assumptions and limitations. A No or NA answer to this question will not be perceived well by the reviewers. 
        \item The claims made should match theoretical and experimental results, and reflect how much the results can be expected to generalize to other settings. 
        \item It is fine to include aspirational goals as motivation as long as it is clear that these goals are not attained by the paper. 
    \end{itemize}

\item {\bf Limitations}
    \item[] Question: Does the paper discuss the limitations of the work performed by the authors?
    \item[] Answer: \answerYes{}
    \item[] Justification: The paper discusses limitations in the Analysis paragraphs, including data scarcity for ViT, resolution constraints for LLM, and the failure of test-time training when the base model is weak. These are elaborated in the discussion and conclusion.
    \item[] Guidelines:
    \begin{itemize}
        \item The answer NA means that the paper has no limitation while the answer No means that the paper has limitations, but those are not discussed in the paper. 
        \item The authors are encouraged to create a separate "Limitations" section in their paper.
        \item The paper should point out any strong assumptions and how robust the results are to violations of these assumptions (e.g., independence assumptions, noiseless settings, model well-specification, asymptotic approximations only holding locally). The authors should reflect on how these assumptions might be violated in practice and what the implications would be.
        \item The authors should reflect on the scope of the claims made, e.g., if the approach was only tested on a few datasets or with a few runs. In general, empirical results often depend on implicit assumptions, which should be articulated.
        \item The authors should reflect on the factors that influence the performance of the approach. For example, a facial recognition algorithm may perform poorly when image resolution is low or images are taken in low lighting. Or a speech-to-text system might not be used reliably to provide closed captions for online lectures because it fails to handle technical jargon.
        \item The authors should discuss the computational efficiency of the proposed algorithms and how they scale with dataset size.
        \item If applicable, the authors should discuss possible limitations of their approach to address problems of privacy and fairness.
        \item While the authors might fear that complete honesty about limitations might be used by reviewers as grounds for rejection, a worse outcome might be that reviewers discover limitations that aren't acknowledged in the paper. The authors should use their best judgment and recognize that individual actions in favor of transparency play an important role in developing norms that preserve the integrity of the community. Reviewers will be specifically instructed to not penalize honesty concerning limitations.
    \end{itemize}

\item {\bf Theory assumptions and proofs}
    \item[] Question: For each theoretical result, does the paper provide the full set of assumptions and a complete (and correct) proof?
    \item[] Answer: \answerNA{}
    \item[] Justification: This paper does not present any new theoretical results; it is an empirical study of existing methods.
    \item[] Guidelines:
    \begin{itemize}
        \item The answer NA means that the paper does not include theoretical results. 
        \item All the theorems, formulas, and proofs in the paper should be numbered and cross-referenced.
        \item All assumptions should be clearly stated or referenced in the statement of any theorems.
        \item The proofs can either appear in the main paper or the supplemental material, but if they appear in the supplemental material, the authors are encouraged to provide a short proof sketch to provide intuition. 
        \item Inversely, any informal proof provided in the core of the paper should be complemented by formal proofs provided in appendix or supplemental material.
        \item Theorems and Lemmas that the proof relies upon should be properly referenced. 
    \end{itemize}

\item {\bf Experimental result reproducibility}
    \item[] Question: Does the paper fully disclose all the information needed to reproduce the main experimental results of the paper to the extent that it affects the main claims and/or conclusions of the paper (regardless of whether the code and data are provided or not)?
    \item[] Answer: \answerYes{}
    \item[] Justification: The paper provides detailed experimental setup: model architectures, hyperparameters (learning rate, epochs, batch size), sampling methods, prompt details, and data splits. It also describes the API used for LLM and the augmentation pipeline. These are sufficient for replication.
    \item[] Guidelines:
    \begin{itemize}
        \item The answer NA means that the paper does not include experiments.
        \item If the paper includes experiments, a No answer to this question will not be perceived well by the reviewers: Making the paper reproducible is important, regardless of whether the code and data are provided or not.
        \item If the contribution is a dataset and/or model, the authors should describe the steps taken to make their results reproducible or verifiable. 
        \item Depending on the contribution, reproducibility can be accomplished in various ways. For example, if the contribution is a novel architecture, describing the architecture fully might suffice, or if the contribution is a specific model and empirical evaluation, it may be necessary to either make it possible for others to replicate the model with the same dataset, or provide access to the model. In general. releasing code and data is often one good way to accomplish this, but reproducibility can also be provided via detailed instructions for how to replicate the results, access to a hosted model (e.g., in the case of a large language model), releasing of a model checkpoint, or other means that are appropriate to the research performed.
        \item While NeurIPS does not require releasing code, the conference does require all submissions to provide some reasonable avenue for reproducibility, which may depend on the nature of the contribution. For example
        \begin{enumerate}
            \item If the contribution is primarily a new algorithm, the paper should make it clear how to reproduce that algorithm.
            \item If the contribution is primarily a new model architecture, the paper should describe the architecture clearly and fully.
            \item If the contribution is a new model (e.g., a large language model), then there should either be a way to access this model for reproducing the results or a way to reproduce the model (e.g., with an open-source dataset or instructions for how to construct the dataset).
            \item We recognize that reproducibility may be tricky in some cases, in which case authors are welcome to describe the particular way they provide for reproducibility. In the case of closed-source models, it may be that access to the model is limited in some way (e.g., to registered users), but it should be possible for other researchers to have some path to reproducing or verifying the results.
        \end{enumerate}
    \end{itemize}

\item {\bf Open access to data and code}
    \item[] Question: Does the paper provide open access to the data and code, with sufficient instructions to faithfully reproduce the main experimental results, as described in supplemental material?
    \item[] Answer: \answerNo{}
    \item[] Justification: The paper does not include a code or data release. However, it provides detailed enough descriptions to allow independent implementation, and the datasets (ARC-AGI) are publicly available.
    \item[] Guidelines:
    \begin{itemize}
        \item The answer NA means that paper does not include experiments requiring code.
        \item Please see the NeurIPS code and data submission guidelines (\url{https://nips.cc/public/guides/CodeSubmissionPolicy}) for more details.
        \item While we encourage the release of code and data, we understand that this might not be possible, so “No” is an acceptable answer. Papers cannot be rejected simply for not including code, unless this is central to the contribution (e.g., for a new open-source benchmark).
        \item The instructions should contain the exact command and environment needed to run to reproduce the results. See the NeurIPS code and data submission guidelines (\url{https://nips.cc/public/guides/CodeSubmissionPolicy}) for more details.
        \item The authors should provide instructions on data access and preparation, including how to access the raw data, preprocessed data, intermediate data, and generated data, etc.
        \item The authors should provide scripts to reproduce all experimental results for the new proposed method and baselines. If only a subset of experiments are reproducible, they should state which ones are omitted from the script and why.
        \item At submission time, to preserve anonymity, the authors should release anonymized versions (if applicable).
        \item Providing as much information as possible in supplemental material (appended to the paper) is recommended, but including URLs to data and code is permitted.
    \end{itemize}

\item {\bf Experimental setting/details}
    \item[] Question: Does the paper specify all the training and test details (e.g., data splits, hyperparameters, how they were chosen, type of optimizer, etc.) necessary to understand the results?
    \item[] Answer: \answerYes{}
    \item[] Justification: Section 5.1 thoroughly describes data splits, hyperparameters (learning rate, epochs, batch size, optimizer), model sizes, and evaluation metrics. The prompt design and API settings are also provided.
    \item[] Guidelines:
    \begin{itemize}
        \item The answer NA means that the paper does not include experiments.
        \item The experimental setting should be presented in the core of the paper to a level of detail that is necessary to appreciate the results and make sense of them.
        \item The full details can be provided either with the code, in appendix, or as supplemental material.
    \end{itemize}

\item {\bf Experiment statistical significance}
    \item[] Question: Does the paper report error bars suitably and correctly defined or other appropriate information about the statistical significance of the experiments?
    \item[] Answer: \answerNo{}
    \item[] Justification: The paper reports single-run accuracy values without error bars because of the computational cost of running multiple seeds on LLM APIs and ViT training. However, the performance gaps are large (46.7\% vs 15.8\%) and the conclusions are robust.
    \item[] Guidelines:
    \begin{itemize}
        \item The answer NA means that the paper does not include experiments.
        \item The authors should answer "Yes" if the results are accompanied by error bars, confidence intervals, or statistical significance tests, at least for the experiments that support the main claims of the paper.
        \item The factors of variability that the error bars are capturing should be clearly stated (for example, train/test split, initialization, random drawing of some parameter, or overall run with given experimental conditions).
        \item The method for calculating the error bars should be explained (closed form formula, call to a library function, bootstrap, etc.)
        \item The assumptions made should be given (e.g., Normally distributed errors).
        \item It should be clear whether the error bar is the standard deviation or the standard error of the mean.
        \item It is OK to report 1-sigma error bars, but one should state it. The authors should preferably report a 2-sigma error bar than state that they have a 96\% CI, if the hypothesis of Normality of errors is not verified.
        \item For asymmetric distributions, the authors should be careful not to show in tables or figures symmetric error bars that would yield results that are out of range (e.g. negative error rates).
        \item If error bars are reported in tables or plots, The authors should explain in the text how they were calculated and reference the corresponding figures or tables in the text.
    \end{itemize}

\item {\bf Experiments compute resources}
    \item[] Question: For each experiment, does the paper provide sufficient information on the computer resources (type of compute workers, memory, time of execution) needed to reproduce the experiments?
    \item[] Answer: \answerYes{}
    \item[] Justification: The paper states that ViT experiments used a 32GB vGPU on a cloud server; LLM experiments used the OpenAI-compatible API. It does not give exact runtime, but the scale is indicated.
    \item[] Guidelines:
    \begin{itemize}
        \item The answer NA means that the paper does not include experiments.
        \item The paper should indicate the type of compute workers CPU or GPU, internal cluster, or cloud provider, including relevant memory and storage.
        \item The paper should provide the amount of compute required for each of the individual experimental runs as well as estimate the total compute. 
        \item The paper should disclose whether the full research project required more compute than the experiments reported in the paper (e.g., preliminary or failed experiments that didn't make it into the paper). 
    \end{itemize}
    
\item {\bf Code of ethics}
    \item[] Question: Does the research conducted in the paper conform, in every respect, with the NeurIPS Code of Ethics \url{https://neurips.cc/public/EthicsGuidelines}?
    \item[] Answer: \answerYes{}
    \item[] Justification: The research uses public benchmark data (ARC-AGI) and does not involve human subjects or sensitive data. There are no ethical concerns.
    \item[] Guidelines:
    \begin{itemize}
        \item The answer NA means that the authors have not reviewed the NeurIPS Code of Ethics.
        \item If the authors answer No, they should explain the special circumstances that require a deviation from the Code of Ethics.
        \item The authors should make sure to preserve anonymity (e.g., if there is a special consideration due to laws or regulations in their jurisdiction).
    \end{itemize}

\item {\bf Broader impacts}
    \item[] Question: Does the paper discuss both potential positive societal impacts and negative societal impacts of the work performed?
    \item[] Answer: \answerYes{}
    \item[] Justification: The paper includes a dedicated "Broader Impact" section that discusses implications for AI capabilities and acknowledges that progress on ARC-AGI could have downstream effects, though no specific negative applications are identified.
    \item[] Guidelines:
    \begin{itemize}
        \item The answer NA means that there is no societal impact of the work performed.
        \item If the authors answer NA or No, they should explain why their work has no societal impact or why the paper does not address societal impact.
        \item Examples of negative societal impacts include potential malicious or unintended uses (e.g., disinformation, generating fake profiles, surveillance), fairness considerations (e.g., deployment of technologies that could make decisions that unfairly impact specific groups), privacy considerations, and security considerations.
        \item The conference expects that many papers will be foundational research and not tied to particular applications, let alone deployments. However, if there is a direct path to any negative applications, the authors should point it out. For example, it is legitimate to point out that an improvement in the quality of generative models could be used to generate deepfakes for disinformation. On the other hand, it is not needed to point out that a generic algorithm for optimizing neural networks could enable people to train models that generate Deepfakes faster.
        \item The authors should consider possible harms that could arise when the technology is being used as intended and functioning correctly, harms that could arise when the technology is being used as intended but gives incorrect results, and harms following from (intentional or unintentional) misuse of the technology.
        \item If there are negative societal impacts, the authors could also discuss possible mitigation strategies (e.g., gated release of models, providing defenses in addition to attacks, mechanisms for monitoring misuse, mechanisms to monitor how a system learns from feedback over time, improving the efficiency and accessibility of ML).
    \end{itemize}
    
\item {\bf Safeguards}
    \item[] Question: Does the paper describe safeguards that have been put in place for responsible release of data or models that have a high risk for misuse (e.g., pretrained language models, image generators, or scraped datasets)?
    \item[] Answer: \answerNA{}
    \item[] Justification: This paper does not release any new models or datasets; it uses existing publicly available benchmarks and LLM APIs. No safeguards are needed.
    \item[] Guidelines:
    \begin{itemize}
        \item The answer NA means that the paper poses no such risks.
        \item Released models that have a high risk for misuse or dual-use should be released with necessary safeguards to allow for controlled use of the model, for example by requiring that users adhere to usage guidelines or restrictions to access the model or implementing safety filters. 
        \item Datasets that have been scraped from the Internet could pose safety risks. The authors should describe how they avoided releasing unsafe images.
        \item We recognize that providing effective safeguards is challenging, and many papers do not require this, but we encourage authors to take this into account and make a best faith effort.
    \end{itemize}

\item {\bf Licenses for existing assets}
    \item[] Question: Are the creators or original owners of assets (e.g., code, data, models), used in the paper, properly credited and are the license and terms of use explicitly mentioned and properly respected?
    \item[] Answer: \answerYes{}
    \item[] Justification: The paper cites all original papers for ARC-AGI, VARC, DeepSeek-R1, and other methods. The ARC-AGI dataset is publicly available under a permissive license; we do not redistribute it.
    \item[] Guidelines:
    \begin{itemize}
        \item The answer NA means that the paper does not use existing assets.
        \item The authors should cite the original paper that produced the code package or dataset.
        \item The authors should state which version of the asset is used and, if possible, include a URL.
        \item The name of the license (e.g., CC-BY 4.0) should be included for each asset.
        \item For scraped data from a particular source (e.g., website), the copyright and terms of service of that source should be provided.
        \item If assets are released, the license, copyright information, and terms of use in the package should be provided. For popular datasets, \url{paperswithcode.com/datasets} has curated licenses for some datasets. Their licensing guide can help determine the license of a dataset.
        \item For existing datasets that are re-packaged, both the original license and the license of the derived asset (if it has changed) should be provided.
        \item If this information is not available online, the authors are encouraged to reach out to the asset's creators.
    \end{itemize}

\item {\bf New assets}
    \item[] Question: Are new assets introduced in the paper well documented and is the documentation provided alongside the assets?
    \item[] Answer: \answerNA{}
    \item[] Justification: The paper does not introduce new assets (datasets, models, or code).
    \item[] Guidelines:
    \begin{itemize}
        \item The answer NA means that the paper does not release new assets.
        \item Researchers should communicate the details of the dataset/code/model as part of their submissions via structured templates. This includes details about training, license, limitations, etc. 
        \item The paper should discuss whether and how consent was obtained from people whose asset is used.
        \item At submission time, remember to anonymize your assets (if applicable). You can either create an anonymized URL or include an anonymized zip file.
    \end{itemize}

\item {\bf Crowdsourcing and research with human subjects}
    \item[] Question: For crowdsourcing experiments and research with human subjects, does the paper include the full text of instructions given to participants and screenshots, if applicable, as well as details about compensation (if any)? 
    \item[] Answer: \answerNA{}
    \item[] Justification: The paper does not involve crowdsourcing or human subjects.
    \item[] Guidelines:
    \begin{itemize}
        \item The answer NA means that the paper does not involve crowdsourcing nor research with human subjects.
        \item Including this information in the supplemental material is fine, but if the main contribution of the paper involves human subjects, then as much detail as possible should be included in the main paper. 
        \item According to the NeurIPS Code of Ethics, workers involved in data collection, curation, or other labor should be paid at least the minimum wage in the country of the data collector. 
    \end{itemize}

\item {\bf Institutional review board (IRB) approvals or equivalent for research with human subjects}
    \item[] Question: Does the paper describe potential risks incurred by study participants, whether such risks were disclosed to the subjects, and whether Institutional Review Board (IRB) approvals (or an equivalent approval/review based on the requirements of your country or institution) were obtained?
    \item[] Answer: \answerNA{}
    \item[] Justification: No human subjects are involved.
    \item[] Guidelines:
    \begin{itemize}
        \item The answer NA means that the paper does not involve crowdsourcing nor research with human subjects.
        \item Depending on the country in which research is conducted, IRB approval (or equivalent) may be required for any human subjects research. If you obtained IRB approval, you should clearly state this in the paper. 
        \item We recognize that the procedures for this may vary significantly between institutions and locations, and we expect authors to adhere to the NeurIPS Code of Ethics and the guidelines for their institution. 
        \item For initial submissions, do not include any information that would break anonymity (if applicable), such as the institution conducting the review.
    \end{itemize}

\item {\bf Declaration of LLM usage}
    \item[] Question: Does the paper describe the usage of LLMs if it is an important, original, or non-standard component of the core methods in this research? 
    \item[] Answer: \answerYes{}
    \item[] Justification: The core methodology uses the DeepSeek-R1 LLM as a reasoning engine. Section 3.1 and 4.1 describe its use, prompting, and output format. The LLM is central to the contribution.
    \item[] Guidelines:
    \begin{itemize}
        \item The answer NA means that the core method development in this research does not involve LLMs as any important, original, or non-standard components.
        \item Please refer to our LLM policy (\url{https://neurips.cc/Conferences/2025/LLM}) for what should or should not be described.
    \end{itemize}

\end{enumerate}

\end{document}